\chapter{Testes, Avalia\c{c}\~{a}o e Discuss\~{a}o}
\label{chap:testes-avaliacao}

\section{Estrat\'{e}gia de Testes}
\label{sec:estrategia-testes}

A avalia\c{c}\~{a}o segue a divis\~{a}o arquitetural apresentada nos cap\'{i}tulos anteriores. A avalia\c{c}\~{a}o concentra-se na contribui\c{c}\~{a}o distribu\'{i}da: contratos, bus, gateway, malha, WebRTC, emparelhamento, permiss\~{o}es e modo de processos. Qualidade de voz, acur\'{a}cia de transcri\c{c}\~{a}o e desempenho de modelos de linguagem pertencem \'{a} base reutilizada do Aurora e exigiriam protocolos pr\'{o}prios.

A inspe\c{c}\~{a}o do reposit\'{o}rio identificou 82 ficheiros de teste, sendo 63 unit\'{a}rios, 18 de integra\c{c}\~{a}o e 1 ficheiro de dados de fixture. O conjunto cobre registro de contratos, BullMQ com Redis, MeshBus, roteamento de malha, failover, emparelhamento, permiss\~{o}es ACL, identidade, permiss\~{o}es de malha, modo de processos, autentica\c{c}\~{a}o WebRTC e ciclo de vida do cliente RTC. Essa organiza\c{c}\~{a}o corresponde aos riscos do projeto: uma falha no contrato quebra descoberta e gateway; uma falha no roteamento envia chamadas ao destino errado; uma falha no controlo de acesso permite uso indevido de capacidades.

A estrat\'{e}gia separa quatro n\'{i}veis. Testes unit\'{a}rios verificam fun\c{c}\~{o}es e classes isoladas, como permiss\~{o}es, identidade, registro de contratos e MeshBus. Testes de integra\c{c}\~{a}o verificam intera\c{c}\~{a}o entre componentes, como roteamento de malha, failover, emparelhamento e permiss\~{o}es. Testes de processo verificam servi\c{c}os separados comunicando por Redis. Testes est\'{a}ticos verificam que os ficheiros-chave ainda compilam como Python. Essa divis\~{a}o cria uma base para afirmar propriedades funcionais; testes de desempenho continuam como etapa separada.

\section{Valida\c{c}\~{a}o Funcional}
\label{sec:validacao-funcional}

A valida\c{c}\~{a}o funcional possui evid\^{e}ncia de desenho, invent\'{a}rio de testes, verifica\c{c}\~{a}o est\'{a}tica e execu\c{c}\~{a}o-alvo de testes unit\'{a}rios no ambiente de revis\~{a}o final desta disserta\c{c}\~{a}o. O comando \texttt{compileall} foi aplicado aos caminhos de implementa\c{c}\~{a}o usados na contribui\c{c}\~{a}o: \verb|app/messaging|, \verb|app/shared/contracts|, \verb|app/shared/auth|, \verb|app/shared/services|, \verb|app/services/gateway|, \verb|app/services/auth| e \verb|scripts/run_processes.py|. O comando terminou com c\'{o}digo 0, indicando que esses ficheiros n\~{a}o possuem erro sint\'{a}tico detectado pelo interpretador dispon\'{i}vel.

Os testes de contratos verificam registro de m\'{o}dulos, decoradores, exporta\c{c}\~{a}o, importa\c{c}\~{a}o, digest e cria\c{c}\~{a}o autom\'{a}tica de m\'{o}dulos. Esses testes sustentam a afirma\c{c}\~{a}o de que os servi\c{c}os podem declarar m\'{e}todos de forma uniforme. Os testes de MeshBus verificam entrega local de eventos, encaminhamento de eventos marcados para malha, bloqueio de reencaminhamento, roteamento local, roteamento remoto, fallback e delega\c{c}\~{a}o de ciclo de vida. Esses casos correspondem ao comportamento descrito na arquitetura de malha \cite{aurora_handoff_doc}.

Os testes de integra\c{c}\~{a}o mapeados verificam cen\'{a}rios de maior escala. O roteamento de malha inclui rota local, rota remota, fallback quando n\~{a}o h\'{a} par, exclus\~{a}o de par obsoleto, sele\c{c}\~{a}o por lat\^{e}ncia e capacidade. O failover inclui timeout remoto, falha de envio, erro de aplica\c{c}\~{a}o, tentativa de segundo par e retorno local ap\'{o}s remo\c{c}\~{a}o de par. O modo de processos inclui arranque de servi\c{c}os, comunica\c{c}\~{a}o, recarga de configura\c{c}\~{a}o, arranque de pilha, isolamento e encerramento. Esses testes dependem de Redis e de execu\c{c}\~{a}o multi-processo; por isso, permanecem fora da execu\c{c}\~{a}o-alvo desta disserta\c{c}\~{a}o.

\section{Avalia\c{c}\~{a}o de Seguran\c{c}a}
\label{sec:avaliacao-seguranca}

A avalia\c{c}\~{a}o de seguran\c{c}a concentra-se no acesso a capacidades, n\~{a}o em uma prova criptogr\'{a}fica completa. O Aurora usa identidades, tokens, escopos, permiss\~{o}es por t\'{o}pico e emparelhamento bilateral para decidir se uma chamada pode atravessar HTTP, WebRTC ou MeshBus. Essa abordagem implementa um controlo de acesso pr\'{o}ximo do modelo RBAC discutido no Cap\'{i}tulo~\ref{chap:enquadramento}, com permiss\~{o}es associadas a contratos concretos \cite{nist_rbac_project,sandhu_ferraiolo_kuhn_2000}.

Os testes de permiss\~{o}es cobrem correspond\^{e}ncia exata, curinga global, curinga por servi\c{c}o, permiss\~{o}es por tipo, interse\c{c}\~{a}o entre permiss\~{o}es de utilizador e escopos de token, identidade an\'{o}nima, identidade de sistema e identidade administrativa. Os testes de WebRTC verificam que pares an\'{o}nimos n\~{a}o podem enviar mensagens arbitr\'{a}rias, que mensagens de autentica\c{c}\~{a}o e chamadas de pairing continuam permitidas, que chamadas n\~{a}o autorizadas recebem bloqueio, e que o timeout encerra pares sem autentica\c{c}\~{a}o. Esses casos correspondem ao desenho do gate de autentica\c{c}\~{a}o WebRTC documentado para o gateway \cite{aurora_gateway_doc}.

O emparelhamento tamb\'{e}m possui testes de fluxo. O conjunto mapeado inclui fluxo completo, c\'{o}digo inv\'{a}lido, troca sem aprova\c{c}\~{a}o, endpoint protegido sem autentica\c{c}\~{a}o e aprova\c{c}\~{a}o sem autoriza\c{c}\~{a}o. Esses casos atacam pontos onde um par poderia obter acesso sem consentimento. O seu alcance permanece limitado a regras de aplica\c{c}\~{a}o; ataques de rede exigem testes de seguran\c{c}a espec\'{i}ficos.

\section{Resultados e Limita\c{c}\~{o}es}
\label{sec:resultados-limitacoes}

A evidência unitária dirigida de 16 de junho de 2026 regista 151 testes aprovados e oito ignorados. É evidência funcional histórica na versão e ambiente arquivados, não validação atual da bateria completa. Os registos funcionais e experimentais posteriores devem ser atribuídos aos seus próprios manifestos, versões de fonte, contagens e datas; nem o número antigo nem uma regressão nova de código comprovam aceitação completa de voz ou distribuição.

O reposit\'{o}rio inclui uma evid\^{e}ncia hist\'{o}rica no ficheiro de handoff: 590 testes aprovados, 40 ignorados e 9 desselecionados por depend\^{e}ncias opcionais ausentes \cite{aurora_handoff_doc}. Esta disserta\c{c}\~{a}o usa esse n\'{u}mero apenas como contexto de estado do reposit\'{o}rio no momento do handoff. A evid\^{e}ncia de junho de 2026 \'{e} menor, mas verific\'{a}vel: ela cobre uma sele\c{c}\~{a}o unit\'{a}ria diretamente ligada aos contratos, gateway, malha e autentica\c{c}\~{a}o. A bateria completa ainda precisa rerodar com Redis e perfis de integra\c{c}\~{a}o.

O retrato de evidência de 15 de outubro de 2026 contém três grupos experimentais distintos. O controlo adicional de pedidos digitados em CPU preserva as suas 900 medições originais no total, com 300 por rota. Os controlos digitados no desktop com GPU contêm 900 medições para cada condição de bus local e HTTP no mesmo desktop, com aquecimentos adicionais excluídos. São medições condicionais da resposta digitada, não medições de áudio completo nem resultados reais entre dispositivos. Um piloto separado processou a gravação fixa integral por reconhecimento de fala real, inferência real de linguagem e síntese real com Piper. Produziu WAV válido e não silencioso, mas falhou a tarefa de referência: uma tentativa, zero resultados corretos, uma falha, zero timeouts e zero amostras medidas de áudio admitidas. O intervalo entre submissão de ficheiro e validação da saída foi 16854,171509 ms; exclui duração da gravação, preparação, encerramento e reprodução. O offload efetivo na GPU não foi estabelecido nesse piloto e seed/threads efetivos não correspondiam aos valores declarados. Uma correção posterior de configuração não altera esses factos históricos. Fila, primeiro token, primeiro áudio e reprodução permanecem fronteiras independentes não medidas. Estes controlos e o piloto falho não sustentam p90/p95 de áudio, overhead exclusivo de protocolo, aceleração causal de hardware, benefício para dispositivo fraco ou conclusão entre regiões. A aceitação individual completa de áudio e superfícies continua a ser pré-requisito das medições entre dispositivos.

\section{Discuss\~{a}o}
\label{sec:discussao}

Os resultados mostram uma contribui\c{c}\~{a}o implementada com testes previstos para os pontos de maior risco e com uma execu\c{c}\~{a}o unit\'{a}ria-alvo bem-sucedida. A limita\c{c}\~{a}o principal permanece no alcance: a disserta\c{c}\~{a}o n\~{a}o converte o invent\'{a}rio completo em evid\^{e}ncia experimental completa porque testes Redis, processos, integra\c{c}\~{a}o ampla, desempenho e redes reais n\~{a}o foram executados. Essa distin\c{c}\~{a}o protege a validade do trabalho.

Mesmo com essa limita\c{c}\~{a}o, a estrutura de testes acompanha os objetivos da proposta. Contratos e registry tratam a modularidade. LocalBus, BullMQBus e MeshBus tratam a transpar\^{e}ncia entre modo local, processo e malha. Roteamento e failover tratam disponibilidade de capacidades. Pairing, tokens, permiss\~{o}es e WebRTC auth tratam controlo de acesso. Processo e Docker tratam separa\c{c}\~{a}o operacional. Essa correspond\^{e}ncia refor\c{c}a que a avalia\c{c}\~{a}o foi desenhada em torno das contribui\c{c}\~{o}es, n\~{a}o em torno de componentes perif\'{e}ricos.

A pr\'{o}xima etapa de valida\c{c}\~{a}o deve transformar o ambiente-alvo criado nesta revis\~{a}o em um protocolo reproduz\'{i}vel de projeto, acrescentando Redis, testes de integra\c{c}\~{a}o, modo de processos, execu\c{c}\~{a}o completa da su\'{i}te e coleta sistem\'{a}tica de logs JUnit. Depois disso, a disserta\c{c}\~{a}o pode atualizar esta sec\c{c}\~{a}o com contagens completas, falhas, tempos de execu\c{c}\~{a}o e limita\c{c}\~{o}es medidas.
